\documentclass[11pt,a4paper]{article}
\usepackage[margin=1.1in]{geometry}
\usepackage{amsmath,amssymb,amsthm,mathtools}
\usepackage[T1]{fontenc}
\usepackage{microtype}
\usepackage{enumitem}
\usepackage{booktabs}
\usepackage[hidelinks]{hyperref}
\hypersetup{pdfauthor={Dritero Mehmetaj},pdftitle={The period-one locus in GL(2,Z)},pdfsubject={11A55}}

\newtheorem{theorem}{Theorem}[section]
\newtheorem{proposition}[theorem]{Proposition}
\newtheorem{lemma}[theorem]{Lemma}
\newtheorem{corollary}[theorem]{Corollary}
\theoremstyle{definition}
\newtheorem{definition}[theorem]{Definition}
\newtheorem{remark}[theorem]{Remark}
\theoremstyle{remark}
\newtheorem{scope}[theorem]{Scope}

\newcommand{\GL}{\operatorname{GL}}
\newcommand{\SL}{\operatorname{SL}}
\newcommand{\Z}{\mathbb{Z}}
\newcommand{\Q}{\mathbb{Q}}
\newcommand{\R}{\mathbb{R}}
\newcommand{\tr}{\operatorname{tr}}
\newcommand{\disc}{\operatorname{disc}}
\newcommand{\lcm}{\operatorname{lcm}}

\title{The period-one locus in $\GL(2,\Z)$:\\
an intrinsic characterization, and a selection that sees only the trace}
\author{Drit\"ero Mehmetaj}
\date{26 August 2026}

\begin{document}
\maketitle

\begin{abstract}
Let $A\in\GL(2,\Z)$ be hyperbolic with non-negative entries. We show that the
continued-fraction expansion of its dominant eigenvalue is purely periodic of period one
\emph{if and only if} $\det A=-1$, in which case $\tr A=m\ge1$ and the eigenvalue is the
$m$-th metallic mean $\lambda_m=\tfrac12(m+\sqrt{m^2+4})$. The locus so described is thus
not a family chosen to make a theorem true: it is cut out by a determinant condition, and
the classification along it is by $(\tr,\det)$ alone.

Two refinements make the statement usable. First, the companion matrix $X_m$ realizes each
attained trace but is not always the unique $\GL(2,\Z)$-class doing so; we compute the
number of primitive $\GL(2,\Z)$-classes of binary quadratic forms of discriminant $m^2+4$
to be $1,1,1,1,1,2,1,1,2,2,1$ for $m=1,\dots,11$, so the class is unique exactly for
$m\le5$ and $m=6$ is the first repetition. Second --- and this is the point of the paper
--- the invariant that singles out the golden member is insensitive to that ambiguity: for
\emph{every} $A$ on the locus with $\tr A=m$, Cayley--Hamilton gives the integer-matrix
identity
\[
   A^2-I\;=\;mA ,
\]
so that $\det(A^2-I)=-m^2$ and, because $A$ is unimodular, the mapping torus of $A^2$ has
first-homology torsion exactly $\Z/m\oplus\Z/m$ --- a knot complement precisely when $m=1$. The selection therefore operates on the whole
period-one locus and not merely on the chosen representatives, which removes the objection a
reader raises first. We also price the period-one restriction against the period-two
relaxation --- there the torsion is $\Z/\gcd(a,b)\oplus\Z/\lcm(a,b)$ and the only knot
complement is again the golden one --- and we state precisely what is \emph{not} obtained:
the passage from a substitution rule to an oriented mapping torus requires typed data that
no result here supplies.
\end{abstract}

\begin{center}\small
\textit{Paper I of four.} The others are II, \emph{A finite spectrum on an infinite lattice} · III, \emph{The boundary-graviton one-loop determinant} · IV, \emph{What a class invariant cannot supply}.
They share a verification convention: every numeral is recomputed by a named script whose
expected output is stated, and every control is designed so that a plausible error makes it fail.
\end{center}

\medskip\noindent\small\textbf{MSC 2020.} Primary 11A55; Secondary 11E16, 11R11, 37B10, 57K32.\\
\textbf{Keywords.} continued fractions, metallic means, binary quadratic forms, mapping torus, figure-eight knot.\normalsize


\section{Introduction}\label{sec:intro}

\subsection{The objection this paper answers}

A theorem of the form ``\emph{inside the family $\mathcal{F}$, the object $x_0$ is the unique
member with property $P$}'' is worth exactly as much as the description of $\mathcal{F}$ is
worth. If $\mathcal{F}$ was assembled by listing the objects that resemble $x_0$, the theorem
is a tautology wearing a quantifier. The sharpest form of the objection is that the family
admits no description independent of its conclusion.

We answer that form of it here for one specific family, and we are explicit about the weaker
forms we do not answer.

The family is the set of hyperbolic $A\in\GL(2,\Z)$ with non-negative entries whose dominant
eigenvalue has a purely periodic continued fraction of period one. This description mentions
no distinguished member. Theorem~\ref{thm:canonical} shows it coincides with the determinant
condition $\det A=-1$; Theorem~\ref{thm:traceonly} shows that the invariant which distinguishes
one member from the rest depends only on the trace, hence is constant on each of the finitely
many $\GL(2,\Z)$-classes sharing a trace, hence cannot be an artifact of which representative
one writes down.

\subsection{Results}

Throughout, $\lambda_m=\tfrac12\bigl(m+\sqrt{m^2+4}\bigr)$ denotes the $m$-th \emph{metallic
mean}, the positive root of $x^2-mx-1=0$; $\lambda_1=\varphi$ is the golden mean and
$\lambda_2=1+\sqrt2$ the silver.

\begin{itemize}[leftmargin=1.4em]
\item \textbf{Theorem~\ref{thm:canonical}} (characterization). Period one $\iff$ $\det A=-1$;
      then $\tr A=m\ge1$, every $m\ge1$ occurs, and the eigenvalue is $\lambda_m$. The
      classification is by $(\tr,\det)$ alone.
\item \textbf{Theorem~\ref{thm:classes}} (how far the representative is canonical). The number
      of primitive $\GL(2,\Z)$-classes at discriminant $m^2+4$ is $1$ for $m\le5$, and $m=6$
      is the first trace carrying two. So $X_m$ \emph{is} the unique class for $m\le5$, and
      the ambiguity is a genuine phenomenon from $m=6$ on rather than a formal caution.
\item \textbf{Theorem~\ref{thm:traceonly}} (the selection sees only the trace). For every $A$
      on the locus with $\tr A=m$ one has $A^2-I=mA$, hence $\det(A^2-I)=-m^2$ and torsion
      $\Z/m\oplus\Z/m$ in the mapping torus of $A^2$, a knot complement iff $m=1$. The identity
      fails off the locus, so it tests the locus rather than an accident of $2\times2$ matrices.
\item \textbf{Proposition~\ref{prop:periodtwo}} (the restriction, priced). Period-two words are
      orientation-preserving already; the period-one family is the diagonal of the period-two
      family; and the only knot complement in the strictly larger two-parameter family is
      still the golden one. Against the one relaxation we tested, the restriction costs
      nothing.
\end{itemize}

\subsection{What is not claimed}\label{sec:notclaimed}

We claim nothing about how one arrives at such an $A$. In particular:

\begin{enumerate}[label=(\alph*),leftmargin=1.8em]
\item \textbf{No canonical passage from a substitution to a $3$-manifold.} Squaring a
      determinant-$(-1)$ incidence matrix produces an orientation-preserving map, but the
      assignment of letters to Dehn twists, the choice of puncture, the orientation, and the
      mapping-torus construction itself are additional typed data. Nothing here supplies them,
      and we do not treat the resulting $3$-manifold as determined by the combinatorics.
\item \textbf{No claim that the seed is forced.} That some minimal-description principle picks
      out this substitution rather than another is not addressed and is not used.
\item \textbf{No uniqueness of the conjugacy class for $m\ge6$}, by Theorem~\ref{thm:classes};
      the results that survive that ambiguity are exactly the ones stated in terms of $(\tr,\det)$.
\item \textbf{No physical interpretation.} None is offered and none is needed for any statement
      in this paper.
\end{enumerate}

\subsection{Relation to what is classical}

Galois' theorem --- a quadratic irrational has a purely periodic continued fraction exactly
when it is \emph{reduced}, i.e.\ $\alpha>1$ and $-1<\bar\alpha<0$ --- is classical
\cite{Perron,HardyWright}. The metallic means and their period-one expansions are elementary
and often observed. The classification of integral binary quadratic forms up to $\GL(2,\Z)$,
and the Latimer--MacDuffee correspondence between conjugacy classes of integral matrices and
ideal classes \cite{LatimerMacDuffee,Taussky}, are likewise standard.

What we believe is not already recorded in this combination is the packaging: the
\emph{equivalence} of the period-one condition with a determinant condition on the locus
(Theorem~\ref{thm:canonical}), the identification of $m=6$ as the exact threshold where the
companion representative stops being canonical (Theorem~\ref{thm:classes}), and above all the
observation that the homological selection is a function of the trace alone
(Theorem~\ref{thm:traceonly}), which is what makes the selection survive the class ambiguity.
We make no priority claim on the classical ingredients.

\section{The characterization}\label{sec:char}

\begin{definition}\label{def:family}
For $m\ge1$ put
\[
   X_m=\begin{pmatrix}m&1\\1&0\end{pmatrix},\qquad
   \phi_m=X_m^{2}=\begin{pmatrix}m^2+1&m\\m&1\end{pmatrix},
\]
so $\det X_m=-1$, $\tr X_m=m$, $\det\phi_m=+1$. Write $M_A$ for the mapping torus of the
once-punctured torus under a homeomorphism inducing $A$ on $H_1$.
\end{definition}

\begin{theorem}[the locus is a determinant condition]\label{thm:canonical}
Let $A\in\GL(2,\Z)$ be hyperbolic with non-negative entries. The continued fraction of its
dominant eigenvalue is purely periodic of period one if and only if $\det A=-1$. In that
case $\tr A=m$ for some $m\ge1$, the eigenvalue is $\lambda_m$, every $m\ge1$ is attained,
and the classification along the locus is by $(\tr,\det)$ alone.
\end{theorem}

\begin{proof}
($\Rightarrow$) Suppose the dominant eigenvalue $x$ has expansion $[m;m,m,\dots]$ of period
one, $m\ge1$. Then $x=m+1/x$, i.e.\ $x^2-mx-1=0$. The companion matrix of $x^2-mx-1$ is
$X_m$, with $\det X_m=-1$ and $\tr X_m=m$. Since $A$ and $X_m$ have the same characteristic
polynomial, $\det A=-1$ and $\tr A=m$.

($\Leftarrow$) Suppose $A$ is hyperbolic with non-negative entries and $\det A=-1$. Its
characteristic polynomial is $x^2-tx-1$ with $t=\tr A$. Non-negativity and hyperbolicity force
$t\ge1$: the entries are non-negative so $t\ge0$, and $t=0$ gives $x^2=1$, not hyperbolic.
The dominant root is $\lambda_t>1$ and satisfies $x=t+1/x$, so its expansion is
$[t;t,t,\dots]$, purely periodic of period one.

For attainment, $X_m$ has non-negative entries, determinant $-1$ and trace $m$ for each
$m\ge1$. For the final clause, the characteristic polynomial of a $2\times2$ matrix is
determined by $(\tr,\det)$, and on this locus $\det$ is constant; so the trace determines the
polynomial, the eigenvalue, and the expansion. Distinct $m$ give distinct traces and hence
distinct conjugacy classes. The converse is not injective on classes --- see
Theorem~\ref{thm:classes} --- only on traces.
\end{proof}

\begin{remark}
The ($\Leftarrow$) direction is Galois' reducedness criterion in disguise:
$\lambda_t>1$ and its conjugate $\bar\lambda_t=\tfrac12(t-\sqrt{t^2+4})$ satisfies
$-1<\bar\lambda_t<0$ precisely because $0<4t$, so $\lambda_t$ is reduced and its expansion is
purely periodic. Period \emph{one} is then the statement that the expansion's single partial
quotient is $\lfloor\lambda_t\rfloor=t$.
\end{remark}

\section{How canonical is the representative?}\label{sec:classes}

Theorem~\ref{thm:canonical} classifies by $(\tr,\det)$, not by conjugacy class. These differ,
and it is worth knowing exactly where.

\begin{scope}[Latimer--MacDuffee]\label{sc:latimer}
Conjugacy classes in $\GL(2,\Z)$ of matrices with a given irreducible characteristic
polynomial correspond to classes of ideals in the associated order
\cite{LatimerMacDuffee,Taussky}. For $x^2-mx-1$ the discriminant is $m^2+4$, and the count is
the primitive $\GL(2,\Z)$-class number of that discriminant. It need not be $1$.
\end{scope}

\begin{theorem}[the threshold is $m=6$]\label{thm:classes}
Let $h^+(m)$ denote the number of $\GL(2,\Z)$-classes of primitive integral binary quadratic
forms of discriminant $m^2+4$. Then
\[
  \bigl(h^+(m)\bigr)_{m=1}^{11}=(1,1,1,1,1,2,1,1,2,2,1).
\]
Consequently $X_m$ is the \emph{unique} class of trace $m$ on the locus for $1\le m\le5$, and
$m=6$ is the first trace at which it is not.
\end{theorem}

\begin{proof}[Computation]
By reduction theory for indefinite forms, each class contains reduced forms $(a,b,c)$ with
$b^2-4ac=D$, $0<b<\sqrt{D}$ and $\sqrt{D}-b<2|a|<\sqrt{D}+b$, and the reduced forms in a
class constitute a single $\rho$-cycle; $\GL(2,\Z)$-equivalence, as opposed to
$\SL(2,\Z)$-equivalence, additionally identifies $(a,b,c)$ with $(-a,b,-c)$. Enumerating
reduced primitive forms and grouping them into $\rho$-cycles modulo that identification gives
the stated counts. The computation is reproduced, with controls, by the script named in
Appendix~\ref{app:verify}.
\end{proof}

\begin{remark}[a recorded disagreement, now resolved]\label{rem:m12}
An independent count communicated to us reported $2$ primitive classes at $m=12$, where our
enumeration gives $3$. \textbf{We have since settled it at $3$}, by re-implementing the reduction
from scratch and validating the implementation against the whole of Theorem~\ref{thm:classes}: the
rebuilt method reproduces $(1,1,1,1,1,2,1,1,2,2,1)$ for $m=1,\dots,11$ exactly before being applied
at $m=12$, where it returns $3$ classes under both $\SL(2,\Z)$- and $\GL(2,\Z)$-equivalence.

\textbf{A caution worth recording, since it cost us a first attempt.} The reduction condition is
$0<b<\sqrt{D}$; implementing it as $b<\lfloor\sqrt{D}\rfloor$ silently discards the case
$b=\lfloor\sqrt{D}\rfloor$, which is legitimate whenever $D$ is not a perfect square. At $D=148$
that single exclusion removes a class and returns $2$. An implementation making this slip agrees
with the correct one at many small $m$ and fails without warning at others, so a count that has not
been validated against a known table should not be trusted.

Nothing in this paper depends on $m\ge12$: the threshold claim concerns the \emph{first}
repetition, at $m=6$.
\end{remark}

\begin{remark}[the trap this exposes]
Restricting to \emph{primitive} forms is not a technicality. Counting all forms, imprimitive
ones included, changes the answer --- content is a $\GL(2,\Z)$-invariant, so imprimitive forms
contribute spurious classes --- and in particular alters the count at $m=4$, which would move
the apparent threshold. The counts above are primitive throughout.
\end{remark}

\section{The selection sees only the trace}\label{sec:traceonly}

This is the section the paper exists for. Theorem~\ref{thm:classes} says the representative is
not canonical past $m=5$. If the invariant that distinguishes the golden member were a function
of the representative, the distinction would be suspect. It is not.

\begin{theorem}[trace-only selection]\label{thm:traceonly}
Let $A\in\GL(2,\Z)$ with $\det A=-1$ and $\tr A=m\ge1$. Then, as an identity of integer
matrices,
\[
   A^2-I\;=\;mA .
\]
Consequently $\det(A^2-I)=m^2\det A=-m^2$ exactly, the mapping torus $M_{A^2}$ has
\[
   H_1(M_{A^2};\Z)\;\cong\;\Z\oplus\Z/m\oplus\Z/m ,
\]
and $M_{A^2}$ has the first homology of a knot complement --- torsion-free $H_1\cong\Z$ ---
if and only if $m=1$. All of this depends on $A$ only through $\tr A$.
\end{theorem}

\begin{proof}
Since $\det A=-1$ and $\tr A=m$, the characteristic polynomial is $\chi_A(x)=x^2-mx-1$, so
Cayley--Hamilton gives $A^2-mA-I=0$, that is, $A^2-I=mA$. Taking determinants,
$\det(A^2-I)=\det(mA)=m^2\det A=-m^2$.

For the homology, the Wang sequence of the mapping torus of a homeomorphism inducing $B$ on
$H_1$ of the fibre gives $H_1(M_B;\Z)\cong\Z\oplus\operatorname{coker}(B-I)$. Here
$B-I=A^2-I=mA$, and $A\in\GL(2,\Z)$ is invertible \emph{over $\Z$}, so $A\Z^2=\Z^2$ and hence
$(mA)\Z^2=m\Z^2$. Therefore
\[
   \operatorname{coker}(A^2-I)\;=\;\Z^2/(mA)\Z^2\;=\;\Z^2/m\Z^2\;\cong\;\Z/m\oplus\Z/m ,
\]
which is trivial precisely when $m=1$. Every quantity named depends only on $m=\tr A$.
\end{proof}

\begin{remark}[why this is the right proof]
The determinant computation alone gives only $|\operatorname{coker}|=m^2$, which leaves the
invariant factors undetermined --- on that information the torsion could a priori be cyclic
$\Z/m^2$. The identity $A^2-I=mA$ settles them at once, exhibiting the cokernel as $\Z^2/m\Z^2$
on the nose. Unimodularity of $A$ is what does the work: it is exactly what permits cancelling
$A$ from $mA$ without changing the quotient.
\end{remark}

\begin{corollary}[the objection is answered]\label{cor:survives}
The selection of $m=1$ is constant on $\GL(2,\Z)$-conjugacy classes of a given trace, and
indeed on the entire set of matrices of that trace on the locus. It therefore does not depend
on which representative of the class one writes down, and in particular is unaffected by the
non-uniqueness of Theorem~\ref{thm:classes} for $m\ge6$.
\end{corollary}

\begin{scope}[what the identity is, and is not, sensitive to]\label{sc:control}
The identity is specific to the locus: for $\det A=+1$ one has $\chi_A(x)=x^2-tx+1$, so
$\chi_A(1)\chi_A(-1)=(2-t)(2+t)=4-t^2$, and $\det(A^2-I)=4-t^2\ne-t^2$. A test that passed on
both loci would be testing arithmetic rather than the hypothesis; this one separates them. The
verification script in Appendix~\ref{app:verify} carries that separation as an explicit control
that must fail off the locus.
\end{scope}

\section{Pricing the period-one restriction}\label{sec:price}

Period one is a restriction. A reader is entitled to ask what it costs. We can answer against
the first relaxation, which is period two; we cannot answer in general, and we say so.

\begin{proposition}[the period-two family]\label{prop:periodtwo}
For $a,b\ge1$ let
\[
   M(a,b)=\begin{pmatrix}a&1\\1&0\end{pmatrix}\begin{pmatrix}b&1\\1&0\end{pmatrix}
         =\begin{pmatrix}ab+1&a\\b&1\end{pmatrix}
\]
be the matrix of a purely periodic continued fraction of period two. Then:
\begin{enumerate}[label=\textup{(\roman*)},leftmargin=2.4em]
\item $\det M(a,b)=+1$ \emph{already}: a period-two word is orientation-preserving without
      squaring;
\item $M(m,m)=\phi_m=X_m^2$, so the period-one family, squared, is exactly the diagonal of the
      period-two family --- which is therefore strictly larger; and
\item $H_1\bigl(M_{M(a,b)};\Z\bigr)\cong\Z\oplus\Z/\gcd(a,b)\oplus\Z/\lcm(a,b)$, so the only
      member of the whole two-parameter family with the homology of a knot complement is
      $(a,b)=(1,1)$ --- the golden one.
\end{enumerate}
\end{proposition}

\begin{proof}
(i) and (ii) are immediate from the displayed product. For (iii),
$M(a,b)-I=\left(\begin{smallmatrix}ab&a\\b&0\end{smallmatrix}\right)$ has entry gcd
$\gcd(a,b)$ and determinant $-ab$, so its invariant factors are $\gcd(a,b)$ and
$ab/\gcd(a,b)=\lcm(a,b)$; apply the Wang sequence as above. The torsion vanishes iff
$\gcd(a,b)=\lcm(a,b)=1$, i.e.\ $a=b=1$.
\end{proof}

\begin{scope}[the price, stated exactly]\label{sc:pricedperiod}
Against the period-two relaxation --- the only one we tested --- the price of restricting to
period one is \emph{zero}: enlarging the family does not produce a competing knot complement.
We have not tested period $\ge3$, and we do not assert that the pattern continues. What we
assert is that the restriction is not doing hidden work at the first place one would look for
it.
\end{scope}

\section{Summary, and the boundary}\label{sec:summary}

The period-one locus is cut out by $\det=-1$ and stratified by trace
(Theorem~\ref{thm:canonical}). The companion representative is canonical exactly up to $m=5$,
with $m=6$ the first repetition (Theorem~\ref{thm:classes}). The homological invariant that
picks out $m=1$ is a function of the trace, hence blind to that ambiguity
(Theorem~\ref{thm:traceonly}, Corollary~\ref{cor:survives}). Relaxing period one to period two
enlarges the family and changes nothing about which member is a knot complement
(Proposition~\ref{prop:periodtwo}).

The boundary is where \S\ref{sec:notclaimed}(a) put it. Everything above is a statement about
matrices, traces and the homology of mapping tori. The step from a substitution rule to a
specific oriented $3$-manifold passes through choices --- letters to Dehn twists, the puncture,
the orientation --- that are typed data supplied from outside. A reader who wants the
$3$-manifold conclusion should regard that step as an assumption, and this paper as settling
only what happens once it is granted.

\appendix

\section{Verification}\label{app:verify}

Every numeral in this paper is recomputed by
\texttt{verify/check\_locus.py}, which runs $15$ checks in four groups and exits non-zero on
any failure. Each group carries at least one \emph{control} --- a check designed so that a
plausible error would make it fail --- because a verification suite in which nothing can fail
verifies nothing.

\begin{center}
\begin{tabular}{@{}llp{6.6cm}@{}}
\toprule
group & discharges & content \\
\midrule
C1 & Thm~\ref{thm:canonical}  & $\lambda_m=m+1/\lambda_m$ to $50$ digits, $m\le40$; $(\tr,\det)(X_m)=(m,-1)$; control that the comparison $\det=+1$ family is live \\
C2 & Thm~\ref{thm:traceonly}  & $A^2-I=mA$, $\det(A^2-I)=-m^2$, and Smith form $(m,m)$ --- verified on all $896$ integer matrices with $\det=-1$ in a bounded box, traces $1$--$26$, \emph{not} only on $X_m$; control that the identity \emph{fails} for $\det=+1$; torsion trivial only at $m=1$ \\
C3 & Thm~\ref{thm:classes}    & primitive class counts at $m^2+4$ for $m\le11$ by independent reduction and $\rho$-cycle enumeration; first repetition at $m=6$; uniqueness for $m\le5$ \\
C4 & Prop~\ref{prop:periodtwo}& $\det M(a,b)=+1$; $M(m,m)$ against the closed form $\phi_m$; torsion $(\gcd,\lcm)$; only $(1,1)$ is a knot complement for $a,b\le24$ \\
\bottomrule
\end{tabular}
\end{center}

The C2 group is the one that matters. It does not check the identity on the family's chosen
representatives; it enumerates every integer matrix of determinant $-1$ in a box and checks all
of them, which is what makes Corollary~\ref{cor:survives} a verified statement rather than a
restatement of the algebra.

\section{Status in the programme register}\label{app:register}

This paper is the first of four drawn from a larger register of questions, each carrying a
status of \texttt{PROVED}, \texttt{REFUTED}, \texttt{CONDITIONAL}, \texttt{EXTERNAL\_BLOCKER}
or \texttt{EMPIRICAL}. We record the register entries this paper bears on, including --- in
fact especially --- the ones it does not close.

\begin{center}
\begin{tabular}{@{}llp{7.2cm}@{}}
\toprule
entry & status & relation to this paper \\
\midrule
\texttt{OA-C0001} & \texttt{REFUTED}    & whether minimal description selects a unique seed independently of encoding. Not addressed here; recorded as refuted in the register. \\
\texttt{OA-C0002} & \texttt{CONDITIONAL}& whether the declared substitution rules select the Fibonacci rule at minimality. Not addressed here. \\
\texttt{OA-C0003} & \texttt{CONDITIONAL}& whether the substitution canonically determines the oriented mapping torus. \S\ref{sec:notclaimed}(a) is this paper's statement of the same limit. \\
\bottomrule
\end{tabular}
\end{center}

\noindent
None of the three is \texttt{PROVED}, and this paper does not change that. What it supplies is
the part of the chain that \emph{is} unconditional --- the characterization of the locus and the
trace-only selection along it --- while leaving the typed-data step exactly as conditional as
the register already records it. We consider stating this more useful than a stronger-sounding
framing that the register would not support.

\begin{thebibliography}{9}
\bibitem{Perron} O.~Perron, \emph{Die Lehre von den Kettenbr\"uchen}, Teubner, 1913.
\bibitem{HardyWright} G.~H.~Hardy and E.~M.~Wright, \emph{An Introduction to the Theory of
  Numbers}, 6th ed., Oxford, 2008, Ch.~X.
\bibitem{LatimerMacDuffee} C.~G.~Latimer and C.~C.~MacDuffee, A correspondence between classes
  of ideals and classes of matrices, \emph{Ann.\ of Math.} \textbf{34} (1933), 313--316.
\bibitem{Taussky} O.~Taussky, On a theorem of Latimer and MacDuffee, \emph{Canad.\ J.\ Math.}
  \textbf{1} (1949), 300--302.
\bibitem{Buell} D.~A.~Buell, \emph{Binary Quadratic Forms}, Springer, 1989.
\bibitem{Thurston} W.~P.~Thurston, \emph{The Geometry and Topology of Three-Manifolds},
  Princeton lecture notes, 1980.
\end{thebibliography}

\end{document}
